
\documentclass[letterpaper,10pt,oneside]{article}
\usepackage[T1]{fontenc}
\usepackage{tgpagella}
\usepackage{setspace}
\usepackage{hyperref}
\usepackage{amsmath}
\usepackage{graphicx}
\usepackage{enumitem}
\graphicspath{ {images/}} %upload your signature to this file
%Change the margins to fit your CV/resume content
\usepackage[left=0.8in, right=0.8in, bottom=0.6in, top=0.7in]{geometry}

%Changes the page numbers - {arabic}=arabic numerals, {gobble}=no page numbers, {roman}=Roman numerals
\pagenumbering{gobble}

%%%%%%%%%%%%%%%%% END OF PREAMBLE %%%%%%%%%%%%%%%%%%%%%

\begin{document}

%%%%%%%%%%%%%%%%% NAME OF APPLICANT %%%%%%%%%%%%%%%%%%%

\noindent  \LARGE{\textbf{Gabriel Pombal Fritsch}} \\
\vspace{-2ex} 
\hrule 
\normalsize

%%%%%%%%%%%%%%%%% CONTACT INFORMATION %%%%%%%%%%%%%%%%%

\vspace{3mm}

\noindent
\begin{tabular}{@{} l @{\hspace{2.2in}} l @{}}
Department of Economics    & Email: gabriel.fritsch@economics.ox.ac.uk \\
University of Oxford      & LinkedIn: \href{https://www.linkedin.com/in/gpfritsch/}{\texttt{linkedin.com/in/gpfritsch}}  \\
Manor Road Building       & Website: \href{https://www.gabrielpfritsch.com/}{\texttt{gabrielpfritsch.com}}  \\
Oxford, OX1 3UQ          & Citizenship: Brazil, Portugal (UK pre-settled status)  \\
\end{tabular}

% EDUCATION
\vspace{3mm}

\noindent \large \textbf{Education}

\vspace{3mm}

\normalsize \noindent Ph.D in Economics, \textbf{University of Oxford} \hfill \begin{minipage}[t]{6cm}\raggedleft
2022 - 2026
\end{minipage}

\vspace{3mm}

\noindent MPhil in Economics, \textbf{University of Oxford} (Distinction) \hfill \begin{minipage}[t]{6cm}\raggedleft
2020 - 2022
\end{minipage}

\vspace{3mm}

\noindent B.A. in Economics, \textbf{Yale University} (Distinction) \hfill \begin{minipage}[t]{6cm}\raggedleft
2013 - 2017
\end{minipage}

% EMPLOYMENT
\vspace{4mm}

\noindent \large \textbf{Employment}
\vspace{3mm}

%\noindent \normalsize \underline{2023} \hspace{9.5mm} Academic Consultant (part-time), Balyasny Asset Management (London, UK)

% \vspace{1mm}

% \noindent \normalsize \underline{2023-Present} \hspace{1mm} Academic Consultant, Balyasny Asset Management (London, UK)

% \vspace{1mm}
\normalsize
\noindent Visiting PhD Student, European Central Bank (Frankfurt, Germany) \hfill \begin{minipage}[t]{4cm}\raggedleft
2025
\end{minipage} \\
\indent \textit{Summer Research Graduate Programme in honour of Ivan Jaccard}

\noindent Academic Consultant, Balyasny Asset Management (London, UK) \hfill \begin{minipage}[t]{4cm}\raggedleft
2023
\end{minipage}

\noindent Economist, Goldman Sachs Global Investment Research (New York, US) \hfill \begin{minipage}[t]{3cm}\raggedleft
2017-2020
\end{minipage}

\noindent Summer Research Analyst, Goldman Sachs Global Investment Research (New York, US) \hfill \begin{minipage}[t]{4cm}\raggedleft
2016
\end{minipage}

% RESEARCH INTERESTS
\vspace{3mm}

\noindent \large \textbf{Research Interests}

\vspace{3mm}

\noindent \normalsize International macroeconomics, monetary economics, macro-development, macroeconometrics

% RESEARCH
\vspace{3mm}

\noindent \large \textbf{Working Papers}

\vspace{3mm} \normalsize

\noindent \textbf{Identifying Fiscal Shocks using Large Language Models} (Job Market Paper)

\noindent \textit{with J. Zachary Mazlish}

\vspace{3mm}

\noindent \textbf{Asymmetric Information and Financial Inclusion: Evidence from Brazil} 

\noindent \textit{with José Renato Haas Ornelas}

\vspace{3mm} \normalsize

\noindent \textbf{Unequal Decoupling} 

\noindent \textit{ with Sergio de Ferra, Kurt Mitman,  and Federica Romei}
\vspace{2mm}

\noindent \small Who gains and who loses from geopolitical fragmentation? We develop a general-equilibrium, quantitative framework that incorporates cross-country heterogeneity in income and wealth inequality, capturing rich income dynamics in both advanced and emerging economies. We then consider two fragmentation scenarios: the closure of the United States or of China to trade and financial flows with the rest of the world. We find that in the United States, only households at the top of the wealth distribution experience welfare gains in both scenarios, while all other households incur losses. Outside the United States, the average welfare outcome is generally positive when the United States closes, as capital becomes more abundant in all countries. The reverse is generally true when China ceases lending. These average effects mask significantly heterogeneous impacts of these shocks on workers and owners of capital in each country.

\vspace{3mm}

\noindent \normalsize \textbf{Sovereign Debt Crises, Fiscal Austerity, and Corporate Default} 

\noindent \textit{with Sergio de Ferra}
\vspace{2mm}

\noindent \small During the Eurozone debt crisis, Italy suffered from an increase in sovereign borrowing costs followed by a sharp fiscal tightening and a reduction in credit to firms. What is the link between sovereign default risk and firms' credit? We propose a transmission channel from sovereign debt crises to firm-level financial frictions through the action of fiscal policy. Using firm-level data for European economies, we show that firm default risk, credit conditions, and performance deteriorate following fiscal consolidations, with a larger impact on firms in the non-tradable sector. We then address this question in a model of endogenous sovereign default in which lending to firms is risky. A fiscal contraction in the crisis country causes a reduction in firms' profits and an increase in their default risk. Firms are heterogeneous in the degree to which the crisis affects them: those in the non-tradable sector are more vulnerable, as the crisis causes demand for their output to fall. Finally, as observed in Italy, the sovereign crisis leads to a contraction in economic activity.

\vspace{3mm}
\newpage
\noindent \normalsize \textbf{Fiscal Policy and Sudden Stops in Open Economies} 
\vspace{1mm}

\noindent \small This paper examines the effects of fiscal policy during balance-of-payments crises and how they differ across countries. Using a quarterly panel dataset of 45 advanced and emerging economies, I find that government spending multipliers on consumption, the trade balance, and the real exchange rate are higher during sudden stops in capital inflows for emerging economies, but not for advanced economies. This non-linearity in spending multipliers is
strongly associated with cross-country differences in the share of foreign currency external debt. To rationalize this, I build a small open economy New Keynesian model with incomplete markets where the state-dependence of fiscal policy is driven by the currency denomination of external debt. When debt is owed in foreign currency -- a salient feature of emerging economies -- a sudden stop generates a debt-deflation mechanism due to currency mismatch. A fiscal expansion that appreciates the real exchange rate therefore relaxes the borrowing constraint and crowds in consumption. However, when households borrow abroad in local currency as in most advanced economies, this mechanism is absent and the difference in fiscal multipliers across financial states disappears.

\vspace{3mm}

\noindent \large \textbf{Honors, Grants and Awards}
\vspace{2mm} \normalsize

\noindent Oxford-OpenAI Research Grant \hfill \begin{minipage}[t]{6cm}\raggedleft
2025 \end{minipage}

\noindent CROSTAG Research Grant \hfill \begin{minipage}[t]{6cm}\raggedleft
2024 \end{minipage}

\noindent Merton Graduate Scholarship in Economics \hfill \begin{minipage}[t]{6cm}\raggedleft
2022 \end{minipage}

\noindent Distinction, MPhil in Economics, University of Oxford \hfill \begin{minipage}[t]{6cm}\raggedleft
2022 \end{minipage}

\noindent Distinction \& Honors in Economics, Yale University \hfill \begin{minipage}[t]{6cm}\raggedleft
2017 \end{minipage}

\noindent Finalist for Best Senior Thesis, Dept. of Economics, Yale University \hfill \begin{minipage}[t]{6cm}\raggedleft 2017 \end{minipage}

\vspace{3mm}

\noindent \large \textbf{Seminars and Conferences} 

\vspace{2mm}\normalsize
\noindent \textbf{2025:} Oxford PhD Seminar, Society for Economic Dynamics Annual Meeting (scheduled)*, European Central Bank (scheduled)

\vspace{0.5mm}
\noindent \textbf{2023-24:} Oxford Macro Away Day, 29th International Conference of Computing in Economics and Finance, Oxford PhD Seminar, Universität Bonn MEF-Seminar*
    
\noindent \textit{* presented by co-author}

\vspace{3mm}

% Academic Experience

\noindent \large \textbf{Academic Experience}
\vspace{2mm}

\noindent \normalsize \underline{Teaching}

\vspace{1mm}
\noindent  Core Macroeconomics (MPhil/PhD, University of Oxford)\hfill \begin{minipage}[t]{6cm}\raggedleft
2023-24 \& 2024-25
\end{minipage}

\noindent  Tools for Macroeconomists (PhD, London School of Economics)\hfill \begin{minipage}[t]{6cm}\raggedleft
Summer 2024
\end{minipage}

\noindent Monetary Economics \& International Finance (Graduate, University of Oxford)\hfill \begin{minipage}[t]{4cm}\raggedleft
Summer 2024
\end{minipage}

\noindent Intermediate Microeconomics (Undergraduate, Yale University) \hfill \begin{minipage}[t]{6cm}\raggedleft
Fall 2016
\end{minipage}

\vspace{3mm}

\noindent \normalsize \underline{Research Assistance}

\vspace{1mm}
\noindent  Sergio de Ferra and Federica Romei (Oxford) \hfill \begin{minipage}[t]{4cm}\raggedleft
2021-2023
\end{minipage}

\noindent François Gerard (QMUL/UCL) \hfill \begin{minipage}[t]{6cm}\raggedleft
2021
\end{minipage}
\vspace{3mm}

\noindent \normalsize \underline{Professional Service}

\vspace{1mm}
\noindent Co-organizer, Macroeconomics Reading Group (Oxford)\hfill\begin{minipage}[t]{4cm}\raggedleft
2023-24, 2024-25
\end{minipage}

% SKILLS
\vspace{4mm}

\noindent \large \textbf{Skills}

\vspace{3mm} \normalsize

\noindent Programming: Python, Matlab, Julia, Dynare, Stata, R

\noindent Languages: English (fluent), Portuguese (native), Spanish (advanced), French (intermediate)

% ADDITIONAL EDUCATION

\vspace{4mm}

\noindent \large \textbf{Additional Education}
\vspace{3mm}

\noindent \normalsize \textbf{Heterogeneous-Agent Macro Workshop}(Goethe Macro Training School) \hfill \begin{minipage}[t]{2cm}\raggedleft
2024
\end{minipage}

\noindent \textit{with Adrien Auclert, Ludwig Straub and Matt Rognlie }

\vspace{3mm}

\noindent \normalsize \textbf{Advanced Tools for Macroeconomists} (London School of Economics) \hfill \begin{minipage}[t]{4cm}\raggedleft
2023
\end{minipage} 

\noindent \textit{with Wouter den Haan \& Pontus Rendahl}

\vspace{3mm}

\noindent \normalsize \textbf{Heterogeneous-Agent Models \& Machine Learning in Macro }(Oxford) \hfill \begin{minipage}[t]{2cm}\raggedleft
2023
\end{minipage}

\noindent \textit{with Jesús Fernández-Villaverde \& Francesco Zanetti}

\vspace{3mm}

\noindent \normalsize \textbf{International Macro \& Trade}(1st Journal of International Economics Summer School) \hfill \begin{minipage}[t]{2cm}\raggedleft
2023
\end{minipage}

\noindent \textit{with Luca Fornaro, Sebnem Kalemli-Ozcan, Enrique Mendoza, Arnaud Costinot, Jonathan Eaton \& Veronica Rappoport} 




\end{document}

